%
% conversation-mode-prd.tex
%
% Product requirements for vox Conversation Mode: a live voice call between
% one human and one Claude Code agent session, mediated by voxd. Written from
% the product-manager altitude --- what the feature must do and for whom,
% not how it is implemented. Use cases and functional requirements are
% traced to each other so a reviewer can check coverage in both directions.
% Architecture, protocol shapes, and the Z specification for the call state
% machine are deliberately out of this document; they follow it, per
% docs/WORKFLOW.md's design-before-implementation gate.
%
% Every use case and requirement is tagged [P1] or [P2]. P1 is the bar to
% hit before anything else is built: an outstanding call with ElevenLabs,
% for a user wearing headphones. P2 broadens that proven experience to
% local providers, Linux, and headphone-free use, against the same
% interfaces P1 defines --- it is sequencing, not a statement that P2 items
% matter less.
%
% Revised after five live spikes (docs/conversation-mode-spikes.md, run
% 2026-08-22) validated or corrected the original requirements against
% real hardware, a real ElevenLabs account, and a real human. Chapter 2
% is new: a technical design chapter synthesizing what the spikes proved
% about how to build this, added at the operator's direction once real
% evidence existed to design from, rather than left for a separate
% document with nothing concrete backing it yet.
%

\documentclass[a4paper,11pt]{report}
\usepackage[margin=1in]{geometry}
\usepackage[colorlinks=true,linkcolor=blue,citecolor=blue,urlcolor=blue]{hyperref}
\usepackage{enumitem}
\usepackage{longtable}
\usepackage{booktabs}
\usepackage{parskip}

\newcommand{\uc}[1]{\textbf{UC-#1}}
\newcommand{\fr}[1]{\textbf{FR-#1}}
\newcommand{\nfr}[1]{\textbf{NFR-#1}}
\newcommand{\pone}{\textsc{[p1]}}
\newcommand{\ptwo}{\textsc{[p2]}}

\begin{document}

\title{Vox Conversation Mode: Product Requirements}
\author{Punt Labs}
\date{August 2026}
\hypersetup{
  pdftitle={Vox Conversation Mode: Product Requirements},
  pdfauthor={Punt Labs},
  pdfsubject={Product Requirements Document}
}
\maketitle
\tableofcontents

\chapter{Product Requirements}

\section{Overview}

Vox today speaks: an agent calls \texttt{mic:unmute} and text becomes audio.
Conversation Mode adds the other half. A human runs \texttt{vox call start}
(or \texttt{/call start} from within a text conversation), talks
naturally, and a Claude Code agent session hears them, thinks, and answers
out loud --- a phone call with an agent, hands and eyes free, while the
human keeps working, walking, or looking away from the screen.

This is not a new product. It is a mode of the existing \texttt{voxd}
daemon: the same process that owns playback and the music program grows the
ability to also own capture, turn-taking, and a live call. Nothing about
Conversation Mode changes how vox speaks when it is not in a call.

The best version of this feature is measured by one thing: whether a human
would rather have this conversation than type it. That bar is not met by
transcription accuracy or voice quality alone --- both existing TTS
providers already clear that bar. It is met by the call \emph{feeling}
alive during the several seconds an agent spends reading files and running
tools, and by an interruption landing as fast as it would with a person on
the other end of a phone.

\textbf{Primary listening scenario: headphones.} The default assumption for
Conversation Mode is a human wearing headphones. This is a product decision,
not an incidental detail: with headphones, the agent's voice never reaches
an open microphone, so the hardest problem in a full-duplex call ---
acoustic echo between the speaker and the mic --- does not exist for the
primary scenario. A secondary scenario, a user on a laptop or desktop
speaker at ordinary office distance without headphones, is real and in
scope, but it is sequenced after the headphone experience is proven, not
solved at the same time.

\subsection{Priorities and phasing}

Conversation Mode ships in two phases, not as one undifferentiated feature:

\begin{description}
  \item[Phase 1 --- prove the bar (\pone).] One provider, ElevenLabs, used
    for both speech recognition and speech synthesis. One listening
    scenario, headphones. One platform, macOS. The goal of Phase 1 is
    narrow on purpose: prove that a call with the best available provider,
    in the easiest audio configuration, is good enough that a person would
    choose to talk instead of type. Everything that makes a call
    \emph{feel} alive --- turn-taking, barge-in, the agent's think-time
    presence, sentence-level response streaming --- is proven here, against
    a single provider, before it is asked to also work across providers or
    platforms.
  \item[Phase 2 --- broaden it (\ptwo).] Local providers on macOS and
    Linux, the no-headphones office scenario, and any additional cloud
    provider. Phase 2 is built against the interfaces Phase 1 already
    defined, not a redesign --- see \fr{21}.
\end{description}

Every use case and functional requirement below carries a \pone{} or
\ptwo{} tag. A \ptwo{} tag is a sequencing decision, not a statement that
the requirement is optional or unimportant.

\subsection{Goals}

\begin{itemize}
  \item \pone{} A human can start, hold, and end a natural spoken
    conversation with one Claude Code agent session, hands-free after
    starting the call, wearing headphones, with ElevenLabs configured.
  \item \pone{} The experience is excellent on that path --- Conversation
    Mode is designed to its best provider first, not down to the weakest
    one it will eventually support.
  \item \ptwo{} A user who wants a fully private conversation --- no audio
    leaving the machine --- can have one, on both macOS and Linux, by
    choosing a local provider configuration.
  \item \ptwo{} The feature works on both macOS and Linux, and without
    headphones, though not necessarily with identical provider options or
    identical audio handling in every configuration.
\end{itemize}

\subsection{Non-Goals}

\begin{itemize}
  \item Multi-party calls (more than one human, or more than one agent).
  \item Calls that span multiple agent sessions or hand off between agents.
  \item Telephony (PSTN, dial-in numbers, SMS).
  \item Video or any non-audio modality.
  \item Provider-side conversational orchestration (e.g.\ ElevenLabs Agents
    Platform or OpenAI's own reasoning model standing in for the Claude Code
    agent). The agent's reasoning is always a Claude Code session; cloud
    providers supply audio transport only.
  \item Automatic mid-call failover between providers, or any other
    runtime resilience behavior. A call uses the provider configuration
    it started with; if that provider becomes unusable, the call ends
    with a clear reason, the same way vox's existing TTS surface fails
    today rather than silently switching providers underneath the user.
    This also covers a provider that degrades rather than fails outright
    (slow, glitchy, briefly unresponsive) mid-call: the system does not
    detect or compensate for degraded-but-still-functioning service ---
    the same scope cut, applied consistently.
  \item A persisted, reviewable record of a call after it ends. Conversation
    Mode optimizes for the live conversation; a post-call transcript is a
    plausible fast-follow, not a launch requirement, and raises its own
    question (does a transcript get written to disk even in the fully
    local, no-network configuration?) that is better answered when it is
    actually being built.
  \item Changing how vox behaves outside an active call.
\end{itemize}

\subsection{Personas}

\begin{description}
  \item[The pair-programmer] Runs Claude Code all day, wants to describe a
    bug or ask a question without stopping to type, especially when their
    hands are occupied or they are away from the keyboard. Wears
    headphones most of the working day already --- the primary persona for
    Phase 1.
  \item[The privacy-conscious user] Wants the option to hold the entire
    conversation on-device, with no audio leaving the machine, matching the
    posture vox already offers via local TTS providers. Phase 2.
  \item[The Linux user] Expects the same feature, accepts that the
    platform-native convenience some macOS users get is not available, and
    should never be told the feature simply does not exist on their
    platform. Phase 2.
\end{description}

\section{Use Cases}

Each use case names the actor, the trigger, the expected flow, and what
``done well'' looks like. Functional requirements in \S\ref{sec:fr} cite the
use cases they satisfy.

\begin{description}

\item[\uc{1} Start a call. \pone{}] \emph{Actor: any user.} The user is
  working in an editor or terminal with an active Claude Code session
  already open, and types \texttt{/call start} into that session. The
  call begins listening
  within a perceptible instant; no setup dialog, no mode selection, is
  required if a working configuration already exists. \emph{Done well:} the
  user does not have to think about which provider is about to answer them,
  and the agent they are now talking to has the context of what they were
  just doing, not a blank start --- see \fr{4}'s 2026-08-23 revision for
  exactly what ``has the context'' means as currently designed (a
  snapshot taken at call start, not a live link to the session's ongoing
  state).

\item[\uc{2} Ask a question and get a spoken answer. \pone{}] \emph{Actor:
  any user.} The user speaks a complete thought, pauses, and the agent
  processes it and answers audibly. \emph{Done well:} the human never
  wonders whether the system heard them, and the reply begins as soon as
  the agent has anything to say --- not only once the agent has finished
  saying all of it.

\item[\uc{3} Interrupt the agent mid-answer. \pone{}] \emph{Actor: any
  user.} The agent is mid-sentence; the human starts talking over it,
  either to correct a wrong assumption or to redirect. \emph{Done well:}
  the agent's voice stops as fast as a human would stop talking when
  interrupted, and what the human said while interrupting is not lost. On
  headphones this does not require echo cancellation to work reliably ---
  the mic only ever hears the human, never the agent's own voice.

\item[\uc{4} Pause mid-thought without ending the turn. \pone{}]
  \emph{Actor: any user.} The human stops speaking briefly --- to think, to
  cough, to react to something in the room --- without intending to hand
  the turn to the agent. \emph{Done well:} the system does not prematurely
  send a half-formed thought, and does not require the human to learn a
  special cadence of speech to avoid this.

\item[\uc{5} Wait through a long agent turn. \pone{}] \emph{Actor: any
  user.} The agent's answer requires reading files, running a search, or
  several tool calls, taking anywhere from a few seconds to tens of
  seconds. \emph{Done well:} the call never goes silent in a way that
  makes the human wonder if it dropped; something --- an acknowledgment, a
  narrated step, an audible cue --- signals the agent is still working.

\item[\uc{6} Hold a fully private conversation. \ptwo{}] \emph{Actor: the
  privacy-conscious user.} The user has configured only local providers, or
  has no network connection at all. \emph{Done well:} the call functions
  end to end with no audio or transcript leaving the machine, on both
  macOS and Linux, even though voice quality and latency may be
  noticeably lower than the cloud experience.

\item[\uc{7} Run the feature on Linux. \ptwo{}] \emph{Actor: the Linux
  user.} The user has no cloud provider configured and no macOS-only local
  option available. \emph{Done well:} the feature still works, using
  whichever local provider is available on that platform, and any setup
  step this requires (e.g.\ a one-time model download) is explained, not
  silently imposed.

\item[\uc{8} Talk without headphones, in an office. \ptwo{}] \emph{Actor:
  any user.} The user is on a laptop or desktop, at ordinary conversational
  distance from the built-in or external mic and speaker --- not across a
  room --- without headphones, in a typical office with some ambient noise.
  \emph{Done well:} the system does not mistake ordinary room noise for
  speech, does not mistake the agent's own voice coming through the
  speaker for a human interruption, does not require silence to function,
  and when it genuinely cannot make out what was said, it asks the human
  to repeat themselves rather than acting on a guess. This use case is
  explicitly not about far-field or across-the-room use --- that is out of
  scope entirely, not merely deferred.

\item[\uc{9} End the call. \pone{}] \emph{Actor: any user.} The user says
  goodbye, runs \texttt{/call stop} (or \texttt{vox call stop}), or simply
  stops engaging. \emph{Done well:} an explicit hangup ends the call
  immediately; an inactive call ends itself after a bounded timeout rather
  than holding the microphone open indefinitely.

\item[\uc{10} Start a call for the first time, unconfigured. \pone{}]
  \emph{Actor: a new user.} No provider has been configured yet and the
  user runs \texttt{/call start}. \emph{Done well:} the user is told what is
  needed to get a call working (an ElevenLabs key, in Phase 1) rather than
  the call silently failing to start or starting into a broken state.

\item[\uc{11} Start a call the configured provider cannot reach. \pone{}]
  \emph{Actor: any user.} The user runs \texttt{/call start} while their
  configured provider is unreachable --- no network, an expired key, a
  provider outage. \emph{Done well:} the call does not start, and the user
  is told why, immediately and in a way they can act on, rather than being
  dropped into a call that will not work.

\end{description}

\section{Functional Requirements}
\label{sec:fr}

Requirements are grouped by capability area. Each cites the use case(s) it
exists to satisfy and carries the same \pone{}/\ptwo{} phase tag.

\subsection{Call lifecycle}

\begin{description}
  \item[\fr{1} \pone{}] The system shall let a user start a call with a
    single action (\texttt{vox call start} or \texttt{/call start}), with
    no required configuration step beyond what already makes vox speak
    today. (\uc{1}) When no provider is configured, the system shall tell
    the user what is needed
    rather than failing silently. (\uc{10})
  \item[\fr{2} \pone{}] The system shall let a user end a call explicitly
    at any time, and shall end an inactive call automatically after a
    bounded timeout. (\uc{9})
  \item[\fr{3} \pone{}] \textbf{Revised.} The system shall communicate
    every call-state transition (ready to listen, agent finished
    speaking, call ending) \emph{audibly}, through the same voice channel
    already used for replies. A visual or text signal is not sufficient
    on its own: live testing (Spike 5) confirmed that a hands/eyes-free
    user genuinely does not see on-screen text, so a state signal that
    depends on being read is not a signal at all for the primary
    scenario. The original wording (``somewhere the user can check it'')
    is too weak and is superseded by this. (\uc{1}, \uc{5})
  \item[\fr{4} \pone{}] A call shall involve exactly one human and one
    Claude Code agent session for the lifetime of that call, and that
    session shall be the user's active session --- carrying whatever task
    context it already had --- not a newly started, context-free session.
    \textbf{Revised (2026-08-23).} The mechanism ratified for \fr{4}
    (Decision 1, \S\ref{sec:design-session-attach}) has changed twice
    since this requirement was written, and the second change narrows
    what ``the user's active session'' can literally mean. The per-turn
    \texttt{claude -p --resume} mechanism (superseded, DES-065) did
    literally resume the same session object; its replacement (DES-066, a
    persistent call agent held alive for the call, spiked and confirmed
    working this session) does not --- it starts a fresh process and is
    handed a \emph{context snapshot} (current task, relevant files,
    recent conversation gist) as its opening message, deliberately
    disconnected from the primary session's own live state so a call
    cannot hang on a primary session that is mid-turn. This still
    satisfies \fr{4}'s underlying intent (the call carries real task
    context, not a blank start) but not its literal wording (the
    \emph{same} session object). \textbf{Operator ratification of this
    reading is required before implementation dispatches} --- see
    \S\ref{sec:design-session-attach}'s updated resolution and
    \nameref{sec:decisions-required} item 1. (\uc{1})
\end{description}

\subsection{Hearing the human}

\begin{description}
  \item[\fr{5} \pone{}] The system shall detect when the human has
    finished a turn, without requiring an explicit ``done'' signal from
    the human. (\uc{2})
  \item[\fr{6} \pone{}] The system shall not end a turn on an ordinary
    thinking pause, a cough, or an incidental noise. (\uc{4}, \uc{8})
  \item[\fr{7} \pone{}] The system shall not treat steady background room
    noise as human speech. (\uc{8})
  \item[\fr{7a} \ptwo{}] Without headphones, the system shall not treat the
    agent's own voice audible through an open speaker as human speech.
    On headphones (\pone{}) this requirement is structurally satisfied ---
    the mic never hears the speaker at all. (\uc{8})
\end{description}

\subsection{Interruption}

\begin{description}
  \item[\fr{8} \pone{}] The system shall let the human interrupt the agent
    while it is speaking, and shall stop the agent's speech promptly on
    interruption. \textbf{Revised:} ``promptly'' is two separate,
    additive latencies (Spike 5), not one number --- see \nfr{1}.
    (\uc{3})
  \item[\fr{9} \pone{}] The system shall not discard what the human said
    while interrupting; that speech shall carry forward as the start of
    the next turn. (\uc{3})
  \item[\fr{10} \ptwo{}] Without headphones, the system shall distinguish
    a genuine interruption from the agent's own voice leaking into the
    microphone, so an interruption is never falsely triggered by the
    system hearing itself. Not required on headphones (\pone{}) for the
    same structural reason as \fr{7a}. (\uc{8})
  \item[\fr{22} \pone{} \textbf{(new)}] The system shall distinguish a
    deliberate interruption from an incidental conversational reaction
    (a brief acknowledgment, exclamation, or backchannel sound like
    ``mm'') by requiring a minimum duration of sustained speech before
    treating audio as an interruption. This holds on headphones too ---
    unlike \fr{10}, it has nothing to do with echo; a human can react
    audibly while wearing headphones. Spike 5 found this by accident (an
    operator's unintentional reaction stopped the agent) and could not
    resolve it from acoustic-duration statistics alone; a working
    threshold (approximately 500\,ms of sustained speech) was reached
    only through direct live tuning against human feedback, and that
    number is one operator's preference in one session, not a validated
    default. See Chapter 2 for the mechanism and \S\ref{sec:oq} for what
    remains open. (\uc{3})
\end{description}

\subsection{The agent's turn}

\begin{description}
  \item[\fr{11} \pone{}] The system shall begin speaking the agent's reply
    as soon as the first complete portion of it is available, rather than
    waiting for the entire reply to be ready. (\uc{2}, \uc{5})
  \item[\fr{12} \pone{}] \textbf{Revised.} While the agent is working on a
    reply, the system shall play a pre-authored, bundled ambience loop
    (quiet workspace foley --- typing, paper shuffle, room tone) on the
    daemon's existing background audio channel, independent of the
    foreground speech channel. The original candidate mechanism ---
    \texttt{mic:music}'s live Music-generation API, ducked under the call
    --- is superseded: Spike 4 found the Music-generation model produces
    repetitive, non-representative content for foley-type sound (it is
    tuned for melody, not realistic ambience) and requires expensive
    per-refresh runtime generation. The corrected mechanism is described
    in Chapter 2. (\uc{5})
  \item[\fr{13} \pone{}] The microphone shall remain available to the
    human throughout the agent's turn, so a human who wants to add or
    correct something is never locked out. (\uc{3}, \uc{5})
\end{description}

\subsection{Providers}

\begin{description}
  \item[\fr{14} \pone{}] The system shall support ElevenLabs as a
    Conversation Mode provider, for both speech recognition and speech
    synthesis, on macOS. (\uc{1})
  \item[\fr{14a} \ptwo{}] The system shall additionally support at least
    one fully local speech provider on each of macOS and Linux. (\uc{6},
    \uc{7})
  \item[\fr{15} \pone{}] The system shall let a user configure which
    provider is used for Conversation Mode, consistent with how providers
    are already configured for vox's existing text-to-speech surface. A
    call uses the configuration it was started with for its entire
    duration. (\uc{1})
  \item[\fr{16} \ptwo{}] The system shall not require any cloud account,
    API key, or network connection to hold a complete call on either
    supported platform, when configured for a local provider. (\uc{6},
    \uc{7})
  \item[\fr{17} \ptwo{}] The system shall clearly account for any setup
    cost unique to a platform's local provider (for example, a one-time
    model download on Linux) rather than presenting it as a silent
    limitation. (\uc{7})
  \item[\fr{18} \pone{}] If the configured provider cannot be reached when
    a call is started, the system shall fail to start the call with a
    clear reason, rather than starting the call and finding out
    mid-conversation. (\uc{11})
  \item[\fr{21} \pone{}] The interfaces between voxd's call logic and a
    speech provider (capture, recognition, turn-taking signals, synthesis)
    shall be defined without assuming ElevenLabs is the only possible
    implementation behind them, so that Phase 2 providers can be added
    later without renegotiating the interface boundary. Phase 1 still
    ships exactly one implementation behind that interface --- the
    abstraction is scoped now so it does not have to be retrofitted, not
    because Phase 1 needs more than one provider to prove its point.
\end{description}

\subsection{Feedback and trust}

\begin{description}
  \item[\fr{19} \pone{}] The system shall never fabricate or guess at what
    the human said when audio is ambiguous or capture fails; it shall ask
    the human to repeat themselves rather than act on a low-confidence
    guess. (\uc{8}, and generally whenever capture confidence is low)
\end{description}

\section{Non-Functional Requirements}

\begin{description}
  \item[\nfr{1} \pone{} (Interruption latency)] \textbf{Revised, validated
    with real data.} Total perceived interruption latency is two
    additive, separately-measured quantities, not one number:
    \begin{enumerate}
      \item \emph{Detection latency} --- sustained speech required before
        the system treats audio as a deliberate interruption (\fr{22}).
        Validated live: 500\,ms was tuned against direct human feedback
        and judged ``good enough'' by the operator. This is a starting
        default, not a proven-optimal or universal constant.
      \item \emph{Termination latency} --- from the stop decision to
        audible silence. Spike 2 measured process-kill time alone
        (0.3--5.3\,ms, real hardware, real interruption) and mistook it
        for the whole quantity; Spike 5 found that killing the playback
        process does not guarantee already-buffered audio stops
        immediately, and this buffer-tail component was never measured.
        \textbf{Open}, not closed --- see \S\ref{sec:oq}.
    \end{enumerate}
    The original single ``under 300\,ms'' target is superseded: it did
    not distinguish these two quantities, and the validated detection
    component alone (500\,ms) already exceeds it. (Supports \fr{8},
    \fr{22}.)
  \item[\nfr{2} \pone{} (Time to first spoken word)] \textbf{Component
    legs confirmed; system-level number now measured too, and it
    falsifies the target.} The 580--683\,ms figure below was, and remains,
    correct for what it measured: the audio-facing legs (STT, TTS,
    detection) with the agent role played by hand by the authoring
    session, per \S\ref{sec:design-session-attach}'s own caveat that this
    was ``provisional for the system as a whole.'' That system-wide number
    is no longer provisional --- Slice 1b (vox-gs9u.2) measured it for
    real, through the actual ratified session-attach mechanism
    (\S\ref{sec:design-session-attach}, Decision 1), on 2026-08-23: a
    20-run benchmark (\texttt{claude -p}, both a resumed real session and
    a fresh session, sonnet and haiku, no accumulated context) put
    time-to-first-word at a \textbf{13--25\,s median}, not
    580--683\,ms --- one to two orders of magnitude over target, dominated
    by process-spawn and hook-bootstrap cost external to the agent's own
    reported turn duration (a fresh-session run reported
    \texttt{duration\_api\_ms: 2243} internally against a 13.2\,s external
    wall clock). \textbf{Per-turn subprocess spawn is confirmed
    unworkable for this NFR and is off the table.} See the resolution
    note in \S\ref{sec:design-session-attach} and \S\ref{sec:known-unknowns}
    for the two-track response: a near-term mitigation (a spoken
    acknowledgment on turn-end, covering the wait, tracked as vox-36xc)
    and a longer-term redesign investigation (a persistent, non-per-turn
    call agent, tracked as vox-hobl). The original component-level
    measurement below is retained as evidence the audio pipeline itself
    is not the bottleneck --- the session-attach mechanism is. (Supports
    \fr{11}.)

    Original component-level measurement, retained for reference:
    measured end to end against a real agent and real tool-use turns
    (Spikes 3 and 5, 5 trials total): 580--683\,ms, flat regardless of
    total reply length (the longest reply in Spike 3 had the lowest
    latency). Spike 1 separately measured the ElevenLabs TTS provider
    floor in isolation, with no agent in the loop --- 145--162\,ms across
    5 trials --- which is a different quantity (provider latency alone)
    supporting the same conclusion from a different angle: it establishes
    the headroom the end-to-end number above has to work with, not an
    additional sample of the same measurement.
  \item[\nfr{3} \ptwo{} (Platform and provider breadth)] Phase 1 targets
    macOS with ElevenLabs only. Phase 2 extends the same interfaces to
    Linux, to local providers on both platforms, and to the no-headphones
    scenario; the platforms and providers are not required to reach that
    experience through identical mechanisms.
  \item[\nfr{4} \ptwo{} (Privacy posture)] In a fully local configuration,
    no audio or transcript content shall leave the machine at any point
    during a call.
  \item[\nfr{5} \pone{} (No regression to existing behavior)] Vox's
    behavior when Conversation Mode is not active shall be unchanged by
    this feature's existence.
  \item[\nfr{6} \pone{} \textbf{(new)} (Audible-only state signaling)]
    No call-state transition (\fr{3}) may depend on a visual or text
    channel as its only signal. A generic, unverified audio cue is not
    sufficient either: Spike 5 found a standalone system chime unreliable
    for a reason never root-caused, while vox's own TTS voice
    (\texttt{mic:unmute}), already proven all session, worked every time.
    The call's own voice channel is the only channel this document treats
    as trustworthy for state signaling until an alternative is separately
    verified.
\end{description}

\section{Open Questions}
\label{sec:oq}

\begin{itemize}
  \item \textbf{Termination latency's buffer-tail component} (\nfr{1}).
    Process-kill time is measured (0.3--5.3\,ms); the time from kill
    signal to actual audible silence, inclusive of any output buffer
    already handed to the OS, is not. This likely requires controlling
    playback through a continuously-owned audio stream rather than
    spawning and killing a player subprocess per segment --- see Chapter
    2, \S\ref{sec:design-termination}.
  \item \textbf{Is 500\,ms the right default for \fr{22}, and should it
    be tunable per user or per session?} Reached through one operator's
    live feedback in one room on one day. Reasonable as a shipped
    default given how it was derived, but not validated across
    speakers, rooms, or microphones.
  \item \textbf{The audible ready-to-listen cue is unresolved.}
    \nfr{6} establishes that it must be audible and that vox's own TTS
    voice is the only channel proven reliable; \fr{3}'s ``ready to
    listen'' transition specifically was never made to work reliably in
    Spike 5 using a non-speech chime. Does the listening cue need to be
    spoken (costing a beat of latency and a few words of vocabulary) or
    is there a non-speech sound that can be independently verified as
    reliable before Phase 1 ships?
  \item \textbf{Multi-session audio contention.} During Spike 5, other
    concurrent Claude Code sessions on the same machine were suspected
    (incorrectly, on investigation) of bleeding audio into the call's
    microphone. The suspicion was wrong in that instance, but the
    underlying question --- what happens to a call when another vox
    client on the same daemon plays audio at the same time --- was never
    actually tested and is not addressed anywhere in this document.
  \item Is a one-time local-model download (\fr{17}) acceptable as part of
    vox's existing install story, or does it need its own onboarding step?
    (Phase 2 question; not blocking Phase 1.)
  \item \fr{18} (fail closed at call start) and the mid-call-degradation
    non-goal are a deliberate simplification to keep scope small --- is
    that the right call, or does early user feedback push toward
    resilience sooner than expected?
  \item How many ambience loops (\fr{12}) ship in the bundle, of what
    length, and does the presence mechanism need a fourth
    \texttt{Format} in the existing Program state machine or a simpler
    static-asset path that bypasses it? Both are architecturally
    supported (Chapter 2, \S\ref{sec:design-presence}); the choice is
    undecided.
\end{itemize}

\chapter{Technical Design Solution}

Chapter 1 states what Conversation Mode must do. This chapter states how,
at the level of components and data flow rather than code --- synthesized
from five live spikes (\texttt{docs/conversation-mode-spikes.md}), not
from advance design. Where a spike left something unresolved, this
chapter says so rather than papering over it; the open items here are the
same items in \S\ref{sec:oq}, restated as design gaps instead of
questions. This chapter is still not the Z specification --- per
\texttt{docs/WORKFLOW.md}'s stateful-audio gate, the call state machine
(\fr{4}'s single-session lifecycle, idle/listening/waiting/speaking) gets
its own formal model before implementation. This chapter is the design
that model will formalize.

\section{Component overview}

A call is one continuously-running loop with five components, each
independently proven by a spike:

\begin{center}
\begin{tabular}{@{}lll@{}}
\toprule
Component & Proven by & Real measurement \\
\midrule
Turn detection (capture $\to$ end-of-turn) & Spikes 2, 5 & see \S\ref{sec:design-turn} \\
Speech recognition (audio $\to$ text) & Spike 1 & 2.2--2.5\,s to first partial \\
Agent turn (text $\to$ reply) & Spike 3, 5 & N/A --- Claude Code itself \\
Sentence-streamed synthesis (reply $\to$ audio) & Spike 3 & 580--683\,ms first word \\
Barge-in (audio $\to$ stop signal) & Spikes 2, 5 (partial) & see \S\ref{sec:design-termination} \\
\bottomrule
\end{tabular}
\end{center}

None of these components block another: turn detection and barge-in
detection are both continuous, audio-only processes that run for the
entire life of a call, independent of which of the other three
components is currently active. \textbf{Correction:} an earlier draft of
this document claimed this meant Spike 5 had ``no orchestration layer to
get wrong.'' That overstates it --- Spike 5's own findings include a
process-coordination bug (finding (a)), an audio-routing bug (finding
(b)), and a barge-in test harness that produced meaningless data because
it never closed the loop between detection and playback (finding (e)),
all genuinely orchestration-level failures between otherwise-correct
components. The corrected claim is narrower: the five components'
\emph{internal} correctness was independently established by different
spikes, which is why fixing each one individually (rather than debugging
an entangled whole) was possible --- not that wiring them together was
uneventful.

\textbf{Design requirement added after Chapter 3's test-plan review:}
these components being independent and continuous means they can call
into the shared call state machine (\S\ref{sec:design-session-attach}
onward) concurrently --- turn detection, barge-in detection, and
sentence-streamed synthesis are three separate processes that may all
observe or attempt to change call state at overlapping moments. If the
state machine is implemented with \texttt{voxd/programs/program.py}'s
\texttt{Program} shape (recommended --- see the design-mission note in
\S\ref{sec:design-termination}), that shape is only safe under
\texttt{Program}'s own documented assumption of a single serializing
writer; \texttt{Program} itself has no internal lock. \textbf{The call
state machine must be reachable only through one serialized dispatch
point} (a queue, an actor, or an explicit lock at the call boundary) ---
this is a requirement on the orchestration layer wrapping the state
machine, not something the pure state class can guarantee by construction.
Where exactly that serialization point lives is part of what
\S\ref{sec:design-session-attach}'s Z specification must settle, not an
implementation detail to discover afterward.

\section{The unresolved core: session attachment}
\label{sec:design-session-attach}

Before any other mechanism in this chapter, one gap has to be named
plainly, because everything else described here is downstream of it and
an earlier draft under-weighted it into a single bullet in Known
Unknowns: \textbf{how a transcribed turn reaches, and how a streamed
reply is pulled from, the user's already-running Claude Code session is
undesigned.} FR-4 requires the call to use the user's active
session --- carrying its existing task context --- not a fresh one. Every
spike that stood in for ``the agent'' actually used \emph{this authoring
session itself}, driven by hand: Spike 3's real agent turns were ``seeded
by a genuinely multi-second real turn in this session,'' and Spike 5's
integration loop ran inside the live authoring conversation with the
operator. No spike attempted programmatic injection of transcribed audio
into an independent, already-running session, nor programmatic extraction
of that session's streamed reply. Turn detection, barge-in, and
sentence-streamed synthesis are all real, all measured, and all
irrelevant if there is no mechanism to actually connect them to a live
agent session that isn't the one authoring this document. \textbf{This
needs at least a sketch of the injection/extraction API shape --- is it a
hook into the Claude Code session surface, a queue the session polls,
does it require the session to run inside \texttt{voxd}'s process space
or communicate over IPC --- before an implementation mission dispatches,}
even if that sketch is only an ADR-style list of rejected and chosen
alternatives, not a full design.

One direct consequence for the rest of this chapter: \textbf{every
latency number reported in \nfr{1} and \nfr{2} was measured with the
agent role played by this authoring session, not through whatever
session-attach mechanism is eventually built.} Those numbers establish
that the audio-facing legs of the pipeline (STT, TTS, detection) are fast
enough --- they do not yet establish that the full system will be, once a
real attach mechanism's own latency is added. Treat \nfr{1}/\nfr{2} as
validated for the components measured, provisional for the system as a
whole.

\paragraph{Resolution (2026-08-23).} The session-attach shape was
designed (the ADR below, Decision 1: a headless \texttt{claude -p
--resume} subprocess per human turn), implemented, and shipped in Slice
1b (vox-gs9u.2). The \emph{provisional} qualifier above is now
\emph{resolved, negatively}: real end-to-end testing (a live operator call
plus a 20-run repeatable benchmark, both documented in \nfr{2}) found the
per-turn subprocess design costs 13--25\,s median per turn, one to two
orders of magnitude over the component-level number this section warned
might not hold. The dominant cost is not the model call --- it is process
spawn plus Claude Code's own SessionStart hook cascade (9--28 hooks
observed per invocation depending on resume-vs-fresh, entirely orthogonal
to this feature's own audio pipeline). \textbf{Spawning a fresh agent
process per conversational turn is confirmed unworkable and is off the
table going forward.} Two tracks follow from this, both tracked outside
this document in beads rather than re-litigated here: a near-term
mitigation to make the current per-turn design tolerable (vox-36xc: a
spoken acknowledgment on turn-end, background audio during the wait, and
a stripped-down harness invocation to shave what latency is honestly
available) and a longer-term architecture investigation (vox-hobl: a
persistent call agent kept alive for the whole call instead of respawned
per turn, running read-only against the codebase/project, with any
codebase-write consequences of the call deferred to a summary the primary
session applies after the call ends). The second track is a genuine fork
from Decision 1 as ratified and needs its own design pass and operator
ratification before implementation, per this project's own workflow
discipline --- it is named here, not designed here.

\paragraph{Further resolution (2026-08-23, same day).} vox-hobl's
architecture investigation ran three narrow spikes, each aimed at killing
the approach quickly if it could not work rather than discovering that a
day into implementation --- \texttt{DESIGN.md} DES-066 has the full
evidence trail; summarized here:

\begin{enumerate}
  \item \textbf{Persistence works.} \texttt{pi --mode rpc --no-session}
    held one process alive across two sequential turns, replying
    correctly to both, with the second turn's own usage report showing it
    reused the first turn's warm context (\texttt{cacheRead}) rather than
    re-establishing it from nothing. The first attempt reproduced a real
    hang (a bundled \texttt{biff-bridge} extension interfering with the
    RPC stdout stream); \texttt{--no-extensions} resolved it outright ---
    root-caused, not routed around.
  \item \textbf{Real tool use holds the latency win.} Re-run with actual
    file-read questions against this repo instead of canned replies: a
    cold turn including a real tool call reached first spoken text at
    6.8\,s and completed at 7.0\,s; a warm follow-up turn completed 4.3\,s
    later. Both an order of magnitude under the 13--25\,s per-turn-spawn
    median DES-065 measured, and that number did zero tool use.
  \item \textbf{Read-only enforcement resolved at the proportionate
    level.} The operator's own scoping narrowed this from ``rock solid''
    to ``the call agent must not collide with repo files a human might be
    editing concurrently through the primary session'' --- it does not
    need to guard against anything touching the primary session's own
    state, since the two are already disconnected by design.
    \texttt{pi}'s native \texttt{--tools read,grep,find,ls} CLI allowlist
    (confirmed against \texttt{pi}'s own documentation, not guessed)
    satisfies this: verified live by prompting the agent to write a probe
    file with only that allowlist set --- no write tool was invoked, and
    the file never appeared on disk.
\end{enumerate}

The mechanism is confirmed workable. \textbf{Further resolved
2026-08-24 (DES-067):} the two smaller pieces this paragraph originally
left open are now decided. The context snapshot is authored by the
\emph{primary} session itself at \texttt{/vox:call start <topic>} --- the
session already holds the task context; handing it off does not need a
retrieval mechanism. Quarry and context7 lookups are folded into the
snapshot the same way, run by the primary session before the call agent
spawns (an initial MCP-in-\texttt{pi} approach was spiked and explicitly
rejected by the operator --- ``do not bloat pi''). The post-call summary
is a structured document (modeled on \texttt{../prfaq/}'s meeting-summary
format: Decisions Made, Action Items, Context Gathered, Beads Touched,
Notes), authored by the call agent itself and read directly by the
still-present primary session when the call ends --- no hook needed.
See DES-067 for the full decision record and rejected alternatives.

What remains before an implementation mission dispatches is the operator
ratification \fr{4}'s revised wording above calls for, plus the
process-supervision choice (a direct \texttt{subprocess.Popen} held by
\texttt{voxd} itself, recommended over \texttt{tmux}/\texttt{pi-tools}'
\texttt{keep} extension to avoid a new runtime dependency for a headless
daemon use case
\texttt{keep} was not built for) --- see DES-066 and
\texttt{docs/conversation-mode-call-agent.tex}'s Z model for the process
lifecycle those pieces attach to.

\section{Decisions required before implementation}
\label{sec:decisions-required}

Three forks in this chapter are explicitly reserved for the operator, per
this project's own workflow rules (design review must surface substantive
issues \emph{before} an implementation mission dispatches, not let them
surface as defects during it). Collected here so none of the three can be
missed by being one bullet among many elsewhere in this chapter:

\begin{enumerate}
  \item \textbf{Session attachment} (\S\ref{sec:design-session-attach}).
    \textbf{Superseded mechanism spike-confirmed; pending operator
    ratification of \fr{4}'s revised wording.} Ratified 2026-08-22 as the
    session-attach ADR's Option D (a headless \texttt{claude -p --resume}
    subprocess per human turn); implemented, shipped in Slice 1b
    (vox-gs9u.2), and confirmed unworkable by real measurement
    (13--25\,s median per turn) on 2026-08-23 --- see
    \S\ref{sec:design-session-attach}'s resolution note. The follow-on
    investigation (vox-hobl: one persistent \texttt{pi --mode rpc} call
    agent held alive for the whole call, not respawned per turn) has now
    been spiked and confirmed working across the three dimensions that
    could have killed it (process persistence, real latency under actual
    tool use, read-only enforcement) --- see the further resolution note
    above and \texttt{DESIGN.md} DES-066. \textbf{What is still open, and
    is an operator decision, not an implementation detail:} this
    mechanism starts the call agent as a fresh process with a context
    \emph{snapshot}, not a literal resume of the primary session, which
    satisfies \fr{4}'s intent but not its original literal wording ---
    \fr{4} above has been revised to name this explicitly and needs
    explicit operator ratification before an implementation mission
    dispatches.
  \item \textbf{Continuous-stream playback ownership}
    (\S\ref{sec:design-termination}): built into production
    \texttt{PlaybackQueue} directly, or as a new sibling component that
    reuses its diagnostics without inheriting its subprocess-per-item
    architecture. The first risks \nfr{5} by modifying code every
    existing vox feature depends on; the second risks the ``two
    pipelines'' duplication this chapter argues against elsewhere. No
    recommendation offered; the trade-off is genuinely close.
  \item \textbf{Presence-signal mechanism}
    (\S\ref{sec:design-presence}): a fourth Program \texttt{Format}, or a
    simpler static-asset path outside the \texttt{Program} state
    machine. Unlike the other two, this chapter's own analysis leans one
    way and should say so rather than leave the trade-off unresolved:
    \textbf{recommend the static-asset path.} The \texttt{Format} path's
    only argument is consistency with the existing state machine, but
    that machine's retry and generation semantics exist to handle a Part
    that can fail to generate --- meaningless for a bundled asset that
    already exists on disk and cannot fail to generate. Consistency with
    machinery that would sit entirely unexercised is not a real argument
    for carrying it. This is a recommendation for the operator to
    confirm or overrule, not a decision already made.
\end{enumerate}

\section{Turn detection}
\label{sec:design-turn}

\textbf{Model: accumulated audible-run duration, not a consecutive-chunk
streak.} The first implementation in Spike 5 tracked consecutive strong
chunks and failed twice: too permissive (2 consecutive chunks, ~60ms) let
transients (coughs, clicks) end a turn falsely; too strict (14
consecutive chunks, ~420ms, an overcorrection) missed genuine continuous
speech, because real speech has brief amplitude dips between syllables
that reset a strict consecutive counter. The corrected model accumulates
audible duration across a \emph{run}, resetting only on a genuine silence
gap ($\geq$200\,ms) --- brief within-word dips do not reset it. This
mirrors Unramble's \texttt{LiveTurnPauseDetector} design (researched
earlier in this project) more closely than the first port did; the first
port's failure was simplifying that design past the point where it still
worked, not a flaw in the underlying approach.

\textbf{Calibration must precede tuning, not substitute for it.} Every
threshold in this system (audible ceiling, strong ceiling, sustain
duration) is relative to an adaptively-tracked noise floor, not an
absolute constant --- but the \emph{multipliers} against that floor were
initially guessed from a prior spike's numbers, in a different room, and
failed silently (returning STT tags like \texttt{[background noise]}
rather than an error) until real RMS data was captured
(\texttt{calibrate.py}: median 0.035, peak 0.377, real silent gaps
mid-utterance) and the model was fixed against evidence. A production
implementation needs an explicit calibration step --- even a few seconds
of ``say something'' at call start --- rather than shipping fixed
multipliers tuned once in one room.

\textbf{Scoping gap: this model is tuned against continuous human
speech's amplitude profile, and it is not established whether that
profile is reliably distinguishable, by duration and gap statistics
alone, from bursty non-speech noise} --- a dog barking intermittently, a
knock, irregular nearby typing --- which by construction can produce the
same ``brief gaps within a longer run'' pattern the model was built to
tolerate. Headphones (P1) remove speaker echo but not ambient room noise
reaching an open mic, so this is not automatically a P2-only
(no-headphones) concern the way \fr{7a}/\fr{10} are; it is unclear from
this design whether P1 also assumes a quiet room, and that assumption
should be stated explicitly rather than left implicit.

\section{Barge-in and its two latencies}
\label{sec:design-termination}

Total perceived interruption latency is \textbf{detection latency +
termination latency}, and they need entirely different engineering:

\textbf{Detection latency} uses the same accumulated-run model as turn
detection (\S\ref{sec:design-turn}), with a longer sustain requirement
(\fr{22}): real interruption behavior sustains for a second or more;
backchannel reactions (``mm,'' a short exclamation) run a few hundred
milliseconds. Spike 5's attempts to derive a clean separating threshold
from acoustic-duration statistics alone failed for three different,
each-diagnosable reasons (methodology flaws in the test, not
inconsistency in human behavior --- see \texttt{conversation-mode-spikes.md}
Spike 5(e) for the detail); direct human-in-the-loop tuning (pick a
value, ask ``too slow or too fast,'' adjust) reached a working number
(500\,ms) where statistical derivation did not. \textbf{Recommendation:
ship this as a live-tunable parameter with a 500\,ms default}, not a
hardcoded constant --- the derivation method that worked once (ask the
user) is also the mechanism that should exist at runtime if the default
turns out wrong for a given voice or room.

\textbf{Termination latency is not process-kill latency.} Spike 5's
prototype stopped playback by sending \texttt{SIGTERM} to an
\texttt{afplay} subprocess per synthesized sentence; Spike 2 measured
that kill taking 0.3--5.3\,ms and treated the interruption problem as
solved. It measured the wrong thing: killing the \emph{process} does not
guarantee audio already handed to the OS output buffer stops
\emph{sounding}, and this buffer-tail latency was never isolated from the
kill-signal latency in any spike. \textbf{Design implication: production
playback should not shell out to a per-segment player subprocess at all.}
It should own a single, continuously-running output audio stream that can
be silenced by clearing its buffer directly, not by killing and
restarting a process. This also removes the per-sentence subprocess-spawn
overhead the spike prototypes paid on every segment.

\textbf{Only the direction is settled here, not the concurrency
discipline around it.} ``Clear the buffer directly'' does not by itself
answer: what happens when a barge-in stop-signal arrives at the same
moment a segment finishes naturally on its own (a race between
turn-completion and interruption)? What happens if a second barge-in
fires while the first stop's buffer-clear is still in flight --- a
stop-during-a-stop? From what thread is the buffer cleared, under what
lock, relative to whatever thread is concurrently writing newly
synthesized audio into that same stream? This document does not need to
answer these at the code level, but per the deferred Z specification
noted at this chapter's start, the single-stream design's ordering
guarantees are a required \emph{input} to that formal model, not a detail
the model can discover while being written.

\textbf{Correction:} an earlier draft of this document described this as
``comparable to \texttt{voxd}'s existing \texttt{PlaybackQueue}.'' On
inspection, \texttt{PlaybackQueue.play\_audio()}
(\texttt{voxd/playback.py:329--466}) spawns one subprocess per file via
\texttt{asyncio.create\_subprocess\_exec}, awaits completion or a
timeout-based kill, and exposes no mid-playback interrupt method at all
--- it is architecturally the same subprocess-per-segment shape this
section is diagnosing as insufficient, not an existing approximation of
the target. \texttt{PlaybackQueue} is the right \emph{place} to add a
continuous-stream capability, so this feature inherits its existing
failure diagnostics instead of duplicating them, but that capability does
not exist in production code today and would be new work.

\textbf{Undecided, and belongs to the operator, not this document:}
whether that new capability is built \emph{into} \texttt{PlaybackQueue}
directly, or as a new sibling component that reuses its diagnostics
without inheriting its subprocess-per-item architecture. The first risks
\nfr{5} (vox's behavior outside a call must be unchanged) by modifying
code every existing vox feature --- chimes, TTS replies, music ---
depends on. The second avoids that regression risk but risks the
``two pipelines instead of one'' duplication this document argues against
in the synthesis section below. This is exactly the kind of fork this
project's own workflow reserves for an explicit operator decision before
an implementation mission dispatches, and it has not been made.

\section{Sentence-streamed synthesis: one pipeline, not two}

Confirmed by Spike 3, built and measured, not merely proposed: vox's
existing batch TTS path (\texttt{core.py:split\_text} $\to$
\texttt{providers/chunked.py:chunked\_synthesize} $\to$
\texttt{stitch\_audio}) and Conversation Mode's incremental path are the
same segmentation problem with one difference --- whether the full input
is known upfront or arrives over time. The design that survived contact
with real code: \textbf{one segmentation core, built on
\texttt{split\_text}'s actual sentence-boundary regex, with two flush
policies on top} --- \texttt{GREEDY} (accumulate to \texttt{max\_chars},
matching today's batch behavior byte-for-byte) for the existing TTS
surface, and \texttt{EAGER} (flush every complete sentence immediately)
for Conversation Mode. The naive port failed once, cleanly diagnosed:
incremental code has to answer ``did the sentence end, or did the input
just end?'' as a real question, which \texttt{split\_text} never has to
ask because it always sees the whole string. That distinction is exactly
what the flush-policy split resolves, without needing two segmenters.

\textbf{Recommendation: build the real implementation directly on this
shape}, not on a re-derivation of it --- the unification was validated in
working code (\texttt{.tmp/spikes/conversation-mode/spike3-streaming/}),
not just argued for.

\textbf{Gap: the async/sync boundary, which Spike 3 itself named as the
dominant unsolved problem, has no design here.} \texttt{voxd}'s handlers
are \texttt{async def} throughout, but \texttt{chunked\_synthesize} and
the per-provider synthesis calls it drives are synchronous, blocking
calls. Spike 3's own write-up states plainly that bridging this boundary,
not the segmenter, is where the real design effort has to go for a
production implementation --- and that bridging design does not exist
yet, here or anywhere else in this document. See \S\ref{sec:known-unknowns}.

\section{Presence signal: a quips-shaped pool, not a single asset and not
\texttt{Program}}
\label{sec:design-presence}

Spike 4 falsified the original candidate (\texttt{mic:music}'s
Music-generation API, ducked under the call) on content grounds ---
repetitive, not representative of real workspace foley, because a
music-generation model is tuned for melody and rhythm, not realistic
ambient sound --- before the level problem (needed to be turned up loud
enough that other audio became uncomfortably loud) was even reached. Two
further rounds of design discussion, after the spike, corrected the
mechanism twice more; this section states where that landed, not the
intermediate steps.

\textbf{The design settled on a fourth existing vox pattern this
document had not previously considered: \texttt{quips.py}.}
\texttt{quips.py} holds fixed-size pools of plain text
(\texttt{STOP\_PHRASES}, \texttt{IDLE\_PHRASES}, and others --- seven to
ten entries each), picked via \texttt{random.choice()} at the moment a
hook fires and passed straight into a live synthesis call
(\texttt{hooks.py:254}). What makes repetition cheap is \texttt{cache.py}:
audio is content-addressed by \texttt{(text, voice, provider)}, so the
first time any phrase plays it is a real API call, and every later
occurrence of the same phrase is a cache hit. This is a fourth pattern
already live in vox, distinct from bundled binary assets (chimes) and
from live dynamic generation with retry/pool semantics
(\texttt{Program}/\texttt{Format}, used for Music), and it turned out to
be the correct fit here --- not either of the two this document
previously weighed against each other:

\begin{itemize}
  \item \textbf{A single bundled loop, replayed every call, was rejected
    on product grounds, not technical ones.} Chimes are fine as static,
    endlessly-repeated assets because they are short and purely
    functional --- nobody notices a chime repeating. A multi-second
    ambient loop replayed identically across every call is exactly the
    kind of content a person does notice, and gets tired of.
  \item \textbf{\texttt{Program}/\texttt{Format} reuse was rejected for
    the reason given in an earlier draft of this section} (retry and
    generation-failure semantics are meaningless for content that is
    authored once and cached, not generated live per call) \textbf{and
    that reasoning still holds} --- but it is now moot rather than a live
    trade-off, because the quips pattern already solves the actual
    problem \texttt{Program}'s pool/rotate/anti-repeat machinery exists
    for (avoiding repetition fatigue) at a fraction of the complexity,
    since nothing here needs to be generated live.
\end{itemize}

The settled design:

\begin{enumerate}
  \item \textbf{Author a fixed pool of Sound Effects prompts} ---
    discussed as approximately five distinct scenarios --- the same
    shape as a \texttt{quips.py} phrase pool, but each entry is a Sound
    Effects prompt (confirmed purpose-built for foley/ambient sound, with
    native seamless looping, via ElevenLabs' Sound Effects API) instead
    of a spoken phrase. \textbf{Correction, carried from an earlier
    draft:} a prior version of this document claimed Sound Effects costs
    ``roughly 1/50th the per-second cost'' of Music; that compared a
    per-second Sound-Effects rate against a per-\emph{track} Music price
    without converting units, and was wrong --- on a like-for-like basis
    Sound Effects is more expensive per second, not cheaper. The real
    economic argument is that each pool entry is synthesized at most once
    per install (see item 2), not the per-second provider rate.
  \item \textbf{Synthesize each pool entry once, cache via the existing
    \texttt{cache.py} infrastructure}, exactly as quip phrases already
    are --- no bundled binary assets shipped in the package, no
    per-call runtime generation. First use of each scenario on a given
    machine pays one real synthesis call; every use after that, of any
    scenario, is a cache hit. Gain-normalization (a live listening test
    found a generated clip needed +6\,dB to be audible at a comfortable
    level) is a one-time authoring-time step per pool entry, the same as
    any other quality check on a quip's phrasing.
  \item \textbf{Select with the same anti-repetition discipline
    \texttt{quips.py} already uses} --- random choice across the pool,
    or a simple last-played exclusion, rather than always the same entry
    or a fixed sequence. This is what actually solves the repetition
    problem chimes don't have and \texttt{Program}'s pool/rotate exists
    for, without needing \texttt{Program}'s generation-failure machinery
    to get it.
  \item \textbf{Runs on the existing background channel, independent of
    speech}, the same channel \texttt{Program}-based Music already uses.
    Confirmed live and by code (\texttt{speech\_handlers.py} has zero
    references to \texttt{Program}, pause, or duck): \texttt{voxd}
    already runs two independent audio channels, foreground
    speech/chime (\texttt{playback.py}) and background Program
    (\texttt{music\_player}). Whether this pool should play through that
    same background channel via a lightweight sibling to
    \texttt{music\_player}, or through a new, smaller mechanism of its
    own, is not decided by this document; either way, no ducking exists
    between foreground speech and background audio today, and whether
    the presence loop needs to duck under speech or can coexist at
    authored levels is untested.
\end{enumerate}

\textbf{Out of scope for this design, raised in discussion but not
committed as a requirement:} a user supplying their own ambience (either
an existing audio file, which fits this same cached/static shape with no
generation step, or a custom prompt generated on demand, which would
need \texttt{Program}/\texttt{Format}'s actual retry-and-generate
machinery, since on-demand generation genuinely can fail transiently the
way a fixed authored pool never does). Chapter 1 has no requirement for
user-customizable presence audio; if that becomes real scope, the
correct design is a hybrid --- quips-style pool for the shipped defaults,
\texttt{Program}/\texttt{Format} reuse specifically for user-triggered
custom generation --- not a reason to revisit the choice made here for
the default case.

\section{Audio-first state signaling}

\nfr{6} is a design constraint, not just a requirement: \textbf{every
call-state cue routes through vox's existing TTS voice channel
(\texttt{mic:unmute}), not a new subsystem.} A standalone system chime
(\texttt{afplay /System/Library/Sounds/Tink.aiff}) was verified working
in isolation but was not reliably heard when invoked from inside the
listener process, for a reason not root-caused in this session ---
plausibly output-device routing differing between a bare shell
invocation and one nested inside a Python subprocess, but not confirmed.
Switching to the same TTS channel already proven reliable all session
worked immediately and consistently. \textbf{Do not build a second,
unverified audio-cue mechanism when a trusted one already exists} ---
if a lighter-weight non-speech cue is wanted later for latency reasons,
it needs its own dedicated reliability verification before anything
depends on it, not an assumption carried over from working once,
standalone, outside the real call path. The chime failure's root cause
should be understood, not just worked around, before it is attempted
again in any form --- it is a plausible symptom of an audio-routing issue
that could resurface in the presence loop's background channel or any
future non-speech cue, not something specific to
\texttt{/System/Library/Sounds/Tink.aiff} alone.

\section{Known unknowns}
\label{sec:known-unknowns}

Carried forward from \S\ref{sec:oq} as concrete design gaps, not
open-ended risk. The first two were missing from this section entirely
until an independent design review caught them; they are the largest
gaps in this chapter, not minor additions.

\begin{itemize}
  \item \textbf{How a transcribed turn reaches, and how a streamed reply
    is pulled from, the user's already-running Claude Code session ---
    designed, implemented, measured, and found too slow as designed.}
    \textbf{Resolved as of 2026-08-23}, superseding the rest of this
    bullet as written during design: the session-attach mechanism was
    designed (session-attach ADR, Decision 1: a headless \texttt{claude -p
    --resume} subprocess per human turn), implemented in Slice 1b
    (vox-gs9u.2), and measured for real. The measurement is the new
    unknown this bullet leaves open: per-turn subprocess spawn costs
    13--25\,s median (20-run benchmark, 2026-08-23), dominated by process
    startup and Claude Code's SessionStart hook cascade, not model
    latency. Per-turn spawn is confirmed off the table; see
    \S\ref{sec:design-session-attach}'s resolution note and \nfr{2} for
    the full data and the two-track response (vox-36xc near-term,
    vox-hobl for the persistent-agent redesign). \textbf{Further resolved
    2026-08-23, same day:} vox-hobl's persistent-agent replacement has
    itself now been spiked and confirmed working (persistence, real
    tool-use latency, read-only enforcement --- see
    \S\ref{sec:design-session-attach}'s further-resolution paragraph and
    \texttt{DESIGN.md} DES-066). \textbf{Further resolved 2026-08-24
    (DES-067):} the context-snapshot content and the post-call summary
    handoff, both named as still undesigned as recently as the previous
    sentence, are now decided --- see
    \S\ref{sec:design-session-attach}'s DES-067 pointer above. What
    remains open is no longer any open design question but the operator
    ratification \fr{4}'s revised wording calls for (the call agent gets
    a context snapshot, not a literal resume of the primary session).
  \item \textbf{Provider-agnostic interfaces (\fr{21}, P1) have no design
    content anywhere in this chapter.} Every spike hardcoded ElevenLabs
    directly --- Spike 1 reused \texttt{providers/elevenlabs.py}
    wholesale; Spike 4 bypassed \texttt{mic:music} to call the
    ElevenLabs SDK directly rather than go through any abstraction.
    \fr{21} requires the capture/recognition/turn-taking/synthesis
    interface boundary to be defined without assuming ElevenLabs is the
    only implementation, and nothing in this chapter sketches what that
    boundary looks like. Silently treating this as solved would be
    wrong; it has not been attempted.
  \item \textbf{The async/sync boundary} between \texttt{voxd}'s
    \texttt{async def} handlers and the synchronous
    \texttt{chunked\_synthesize}/provider synthesis calls they currently
    drive --- named by Spike 3 itself as the dominant unsolved
    engineering problem, ahead of the segmenter --- has no design here.
  \item \textbf{Termination's buffer-tail latency} is unmeasured. The
    architectural fix (\S\ref{sec:design-termination}, a continuously-owned
    output stream instead of per-segment subprocess playback) is proposed
    but not built or measured.
  \item \textbf{Whether the new continuous-stream playback capability
    modifies production \texttt{PlaybackQueue} directly or is built as a
    separate component} is an explicit operator decision, not yet made
    (\S\ref{sec:design-termination}).
  \item \textbf{The 500\,ms interruption-detection default} is one
    operator's live-tuned preference, not validated across voices, rooms,
    or microphones. Ship it as a runtime-tunable parameter, not a
    constant.
  \item \textbf{Multi-session audio contention} --- what a call does when
    another vox client on the same daemon plays audio concurrently --- was
    never tested. Spike 5 suspected it incorrectly in one instance; the
    underlying scenario is untested, not disproven.
  \item \textbf{Presence-signal ducking against speech} is untested; the
    two channels are confirmed independent today, and whether that is
    sufficient or needs explicit ducking is open.
  \item \textbf{Ambience-Format vs.\ static-asset path} for
    \S\ref{sec:design-presence} item 4 is an implementation choice not
    yet made.
  \item \textbf{FR-19 (never fabricate on ambiguous audio)} has no design
    coverage here, despite two directly relevant spike findings: Spike
    1's ElevenLabs partial transcripts self-correcting mid-stream, and
    Spike 5's turn detector producing false triggers before it was
    fixed. How the production system surfaces low-confidence capture to
    the human is undesigned.
  \item \textbf{Mid-call self-correction.} Spike 3 found that a live
    agent composing a reply sentence-by-sentence will, at least
    occasionally, speak a first sentence that a later sentence's
    reasoning contradicts or walks back --- and flagged this explicitly
    as belonging in the design mission's open questions. Whether the
    agent should audibly self-correct mid-call the way a human would is
    not addressed anywhere in this document.
  \item \textbf{Which detector governs a human speaking during the
    agent's silent think-time (the waiting state, before any reply audio
    starts)} is not stated. Barge-in detection is described as active
    ``while it is speaking''; turn detection covers human speech
    generally; neither section says which one applies, if either,
    between a transcribed turn being sent and the reply's first audio.
    This is exactly the kind of state-machine edge case the deferred Z
    specification needs to resolve, and it should be named before that
    modeling starts, not discovered during it.
  \item \textbf{A stuck or free-running detector has no described
    recovery.} Both turn detection and barge-in detection are described
    as continuous processes running the whole life of a call, but nothing
    addresses a detector thread that hangs (the call goes silently deaf;
    \uc{9}'s bounded timeout is a call-level backstop, not a
    component-health check) or one that free-runs and never accumulates
    silence (the human's turn never ends). Per \nfr{6}, the TTS voice is
    the only trusted signaling channel --- but a stuck detector by
    definition produces no signal to send through it, so the human would
    have no way to know this is happening mid-call.
  \item \textbf{A synthesis call that never returns mid-turn} is
    unaddressed. \fr{18} covers a provider unreachable \emph{before} a
    call starts; nothing covers a provider that accepts the call and then
    hangs on a synthesis request partway through. Likely a one-sentence
    fix (a timeout on the synthesis call, ending the turn the same way
    \fr{18} ends the call) rather than new design, but it is not stated
    anywhere as written.
  \item \textbf{No debugging story for a live call.} A user-facing
    transcript is an explicit Non-Goal, correctly, for launch scope --- but
    an internal, developer-facing record of state transitions, detector
    decisions, and per-segment latencies is a different thing and is not
    mentioned. Given Phase 1's entire purpose is validating a subjective
    bar (``would a human rather have this conversation than type it''),
    there is currently no way to determine \emph{why} a given call felt
    bad --- a false barge-in trigger, a slow synthesis segment, a missed
    turn end --- after the fact.
  \item \textbf{Whether anything in this design is actually macOS-specific
    is not established.} P1 is scoped to macOS, framed as a complexity
    reduction alongside headphones-only and single-provider. Headphones
    genuinely eliminate an entire class of problem (echo cancellation);
    single-provider genuinely defers multi-provider complexity. But
    nothing in this chapter identifies what part of the design --- if
    any --- is actually coupled to macOS versus incidentally macOS
    because that is where the spikes happened to run. If the design is
    in fact platform-portable, ``macOS-only for P1'' is not reducing
    complexity, it is leaving Linux untested, which is a different
    (and smaller) claim than the phasing rationale in Chapter 1 makes.
\end{itemize}

\chapter{Test Plan}

Chapter 2 designs the system; this chapter designs how it is proven
correct, including what has to be \emph{built} to make that possible ---
this project has no existing test tier for continuous audio-hardware
behavior, and two of the new components have no fake/test-double
equivalent of the ones already established elsewhere in vox. Grounded in
real, present conventions: the four-tier testing pyramid and Humble
Object pattern already documented for this project (\texttt{TESTING.md},
\texttt{punt-kit/standards/python.md}), and the structural, non-inheritance
fake style already in use (\texttt{tests/\_program\_fakes.py}'s
\texttt{FakeProgramGateway}, which records every call and lets a test
control results without a live daemon).

\section{Testing tiers}

\begin{center}
\begin{tabular}{@{}p{2.3cm}p{5.6cm}p{3.3cm}p{2.5cm}@{}}
\toprule
Tier & What it tests & Status for this feature & Runs in CI \\
\midrule
1. Unit & Pure logic: turn/barge-in detectors, the call state machine,
  segmenter flush policies, presence-pool selection & \textbf{Needs new
  fakes} (\S\ref{sec:test-di}) & Yes \\
2. Integration & Cross-component wiring: provider protocol conformance,
  the commands-layer surface, the Z spec's transitions exercised
  end-to-end in memory & \textbf{Needs a fake session-attach double}
  (\S\ref{sec:test-di}) & Yes \\
3. Subprocess/E2E & Wire protocol, real provider calls, real (non-human)
  process lifecycle & \textbf{Exists in prototype form only} --- every
  spike script (\texttt{.tmp/spikes/conversation-mode/*}) is this tier,
  unformalized & Optional \\
4. SDK & End-to-end with a real Claude Code session, costs money &
  \textbf{New for this feature} --- prior vox SDK-tier tests never
  attached to an independent live session; this is Slice 1b's core
  claim and has no existing harness & No \\
5. \textbf{Live audio, human-judged (new tier)} & Does it \emph{sound}
  right --- barge-in feel, presence-signal naturalness, overall call
  quality & \textbf{Does not exist as a tier today}; this project's
  audio-demo gate (\texttt{docs/WORKFLOW.md}) is the closest precedent
  but has never been formalized into a repeatable checklist for a
  \emph{continuous, full-duplex} feature & No, by
  construction \\
\bottomrule
\end{tabular}
\end{center}

\textbf{Tier 5 must be built, not assumed.} The project's existing
audio-demo discipline (\texttt{docs/WORKFLOW.md}) asks a human to confirm
audio sounds right, but every prior use of it was for discrete,
single-shot playback (a chime, a spoken reply) with an obvious pass/fail.
Conversation Mode's audio-demo gate covers a \emph{continuous, two-way}
experience --- barge-in timing, presence-signal fatigue, whether a call
feels alive during a long wait --- which needs its own checklist, not a
reuse of the existing one by name only. This chapter treats writing that
checklist as a required Tier 5 deliverable, owed by Slice 5/Slice 6
(where failure modes and call endings are defined) rather than left
implicit.

\textbf{Tier 3 exists only as five throwaway spike scripts.} Each of
Spikes 1, 2, 3, and 5 built a real, working subprocess-level test
harness against real hardware and real ElevenLabs calls
(\texttt{.tmp/spikes/conversation-mode/*}), explicitly marked throwaway
per \texttt{docs/conversation-mode-spikes.md}. \textbf{Recommendation:
promote the patterns these scripts already validated} (measuring
sentence-1-ready-to-first-audio-byte, measuring detection-to-kill
latency, driving a real turn through a real session) into permanent,
optional-in-CI Tier 3 tests, rather than re-deriving the same harnesses
from scratch during implementation. The scripts are throwaway as
\emph{code}; the measurement patterns they proved out are not.

\section{Design requirements for testability}
\label{sec:test-di}

Each requirement below exists because a component in Chapter 2 currently
has no seam for a test double, and the fix is the same one vox's
\texttt{FakeProgramGateway} and its five-provider \texttt{TTSProvider}
already demonstrate: a structural (Protocol-shaped, non-inheritance)
interface, injected, not owned, by the component under test.

\begin{enumerate}
  \item \textbf{Detectors take PCM chunks as plain method arguments;
    they never own an audio device.} The turn detector and barge-in
    detector (\S\ref{sec:design-turn}, \S\ref{sec:design-termination})
    must be constructible and fully testable with zero real microphone
    access --- feed a chunk, assert the signal. This is not a new
    invention: it is what every spike script already did by hand
    (\texttt{calibrate.py} fed captured PCM into pure functions). The
    requirement is to keep that shape in production code, not let the
    detector's constructor reach for \texttt{sounddevice} directly the
    way the spike prototypes did for expedience.
  \item \textbf{Chunk duration/timestamp is passed as data, never read
    from the wall clock.} \textbf{Correction, from peer review:} the
    accumulated-run model (\S\ref{sec:design-turn}) depends on real-time
    durations --- a $\geq$200\,ms silence gap resets a run, ${\sim}$500\,ms
    sustained speech triggers an interruption. If the detector reads
    \texttt{time.monotonic()} (or equivalent) internally rather than
    taking elapsed time as a parameter, every test built against item 1
    becomes wall-clock-dependent and flaky --- the exact failure mode
    item 1 already polices for the microphone, missed for the clock. The
    detector's public interface must accept duration/timestamp
    explicitly.
  \item \textbf{An \texttt{STTProvider} protocol, following the same
    \emph{pattern} as \texttt{TTSProvider} --- not the same shape.}
    \textbf{Correction, from peer review:} an earlier draft of this
    chapter said ``conforming to the same shape as the existing
    \texttt{TTSProvider}.'' Read literally that is wrong:
    \texttt{TTSProvider}'s members (\texttt{name}, \texttt{default\_voice},
    \texttt{supports\_expressive\_tags}, voice catalog/resolution
    methods) are half voice-catalog logic with no STT analogue, and
    \texttt{STTProvider} needs partial/streaming transcript delivery
    (Spike 1's self-correcting mid-stream partials, Chapter 2) and a
    confidence signal (\fr{19}: ask the human to repeat rather than act
    on a low-confidence guess) that \texttt{TTSProvider} has no
    counterpart for. The pattern that transfers is: \texttt{@runtime\_checkable
    Protocol}, a \texttt{check\_health() -> list[HealthCheck]} method,
    one conforming implementation per vendor, structural (not
    inheritance-based) conformance --- not the member list. A fake STT
    provider built against this returns a scripted transcript
    \emph{and} a scripted confidence value without a network call, the
    same way any \texttt{TTSProvider} test substitutes a fake without
    hitting ElevenLabs.
  \item \textbf{A session-attach protocol, with a fake session double
    that scripts a \emph{sequence} of reply chunks, not one reply} ---
    the single largest gap this chapter adds to Chapter 2.
    \textbf{Correction, from peer review:} an earlier draft specified the
    fake as returning one scripted reply value, which cannot exercise
    \fr{11} (speaking begins on the reply's first complete portion, not
    once the whole reply exists) --- the exact behavior that makes a
    long agent turn (\uc{5}) feel alive. The fake must return a scripted
    sequence of reply chunks with independently controllable timing
    between them, mirroring what it stands in for: a stream, not a
    call-response pair. Structurally the same discipline as
    \texttt{FakeProgramGateway} (records every injected transcript,
    caller controls what comes back) --- but the ``what comes back''
    shape has to be a stream to be useful. \textbf{This fake does not
    exist today and is new required work.}
  \item \textbf{A playback-sink protocol, with an in-memory fake ---
    necessary but not sufficient on its own.} \textbf{Correction, from
    peer review:} an earlier draft implied a call-recording fake alone
    makes the buffer-clear ordering questions
    (\S\ref{sec:design-termination}: stop-during-a-stop, buffer-clear
    racing a concurrent writer) unit-testable. It does not --- those are
    genuinely concurrent races that only exist when two call sites
    contend, and a fake driven single-threadedly proves nothing about a
    race. The in-memory sink is still required (it removes the
    real-audio-output dependency), but it must be paired with a
    concurrency/interleaving harness --- controlled \texttt{asyncio} task
    scheduling with deterministic yield points, or a threaded stress
    test run N times --- that actually drives the race, not just a fake
    that records calls made from one thread.
  \item \textbf{The call state machine is a pure class plus an injected
    policy, reachable only through a single serialized dispatch point.}
    \textbf{Correction, from peer review:} mirroring
    \texttt{voxd/programs/program.py}'s \texttt{Program} for the pure
    state/transition shape is right, but an earlier draft said to mirror
    it ``exactly'' without also copying the discipline that makes it
    safe: \texttt{Program} has no internal lock and does bare
    \texttt{self.\_state = self.\_state.with\_updates(...)} assignments,
    safe only because \texttt{Program}'s own docstring states something
    upstream serializes every call into one writer. The call state
    machine's actual callers --- turn detection, barge-in detection, and
    sentence-streamed synthesis --- are three independently-running,
    continuous processes (\S``Component overview'') that can race
    exactly the way \S\ref{sec:design-termination} already flags as
    unresolved. Copying \texttt{Program}'s pure-class shape without also
    requiring a single serialized dispatch point (a queue, an actor, or
    an explicit lock at the call boundary) reproduces a data race, not a
    design --- this is a requirement on Chapter 2's orchestration layer,
    not something the pure class or its fake can supply alone.
  \item \textbf{Presence-pool selection takes an injectable source of
    randomness / play-history}, not a bare call to \texttt{random.choice}
    at the call site. A test needs to assert anti-repeat behavior
    deterministically (``the same entry never plays twice in a row''),
    which requires controlling what ``last played'' and ``next pick''
    resolve to, the same way any seeded-random test would.
  \item \textbf{The async/sync bridge between \texttt{voxd}'s
    \texttt{async def} handlers and the synchronous STT/TTS provider
    calls must itself be injectable.} \textbf{Added per peer review:} an
    earlier draft of this chapter had no testability treatment at all
    for the boundary Chapter 2 names as ``the dominant unsolved
    engineering problem'' (\S``Sentence-streamed synthesis''), despite
    applying exactly this discipline to the microphone (item 1) and the
    clock (item 2). Whatever bridges the boundary --- a thread pool
    executor, subprocess isolation, or an async-native provider client
    --- must be substitutable with (a) a fake that completes
    synchronously, proving the bridge introduces no deadlock, and (b) a
    fake that stalls, proving a slow synthesis call does not block the
    event loop for its duration.
\end{enumerate}

\section{Test coverage requirements}

\textbf{Corrected per peer review: items 1, 2, 3, and 5 below originally
stated intent with no enforcement mechanism} --- compare to item 4, which
was already enforceable the way \texttt{PL-TT-2} is elsewhere in this
project's standards (a tool check, not a sentence). Each now names a
concrete, greppable artifact.

\begin{enumerate}
  \item \textbf{Every operation in the deferred Z specification has at
    least one Tier 1 test exercising it}, once that spec exists (Slice
    1a). \textbf{Enforcement:} test names include the Z-spec operation
    name verbatim (per \texttt{docs/WORKFLOW.md}'s ``assert the modeled
    properties by name''); a script diffs the operation names in the
    \texttt{.tex} Z spec against test names under
    \texttt{tests/conversation\_mode/} and fails on any operation with
    zero matches. A state transition with no corresponding test is
    exactly the class of defect the Z-spec gate exists to prevent.
  \item \textbf{Every functional and non-functional requirement in
    Chapter 1 maps to at least one test at a named tier.}
    \textbf{Enforcement:} a coverage-matrix file (FR/NFR id $\to$ tier
    $\to$ test path) checked in alongside the tests; a script diffs its
    FR/NFR ids against this document's \fr{} and \nfr{} labels and fails
    on any requirement missing a row. Where a requirement is genuinely
    only checkable by ear (NFR-6's ``does this sound right,'' FR-12's
    presence-signal naturalness), the mapped tier is 5, not a missing
    row. \fr{19} is the counter-example worth naming explicitly: despite
    sounding like a judgment call, it is Tier-1-testable and cheap once
    the \texttt{STTProvider} fake (\S\ref{sec:test-di} item 3) exists ---
    script a low-confidence result, assert the system asks the human to
    repeat rather than acting on it. Not everything audio-adjacent
    defaults to Tier 5.
  \item \textbf{Every Known Unknown that gets resolved during
    implementation ships with a regression-style test or measurement
    harness, not just a one-time manual check.} \textbf{Enforcement:}
    Tier 3 regression harnesses live under a required, named directory
    (\texttt{tests/conversation\_mode/harnesses/}) so their presence is
    greppable, not merely encouraged. Termination's buffer-tail latency
    (currently unmeasured) and \textbf{multi-session audio contention}
    (Chapter 2's Known Unknowns; untested by any spike, and the thing
    Spike 5 incorrectly suspected mid-session) are this item's two
    named worked examples --- once either is resolved, its measurement
    becomes a harness in that directory, re-runnable, not a number
    recorded once and never checked again.
  \item \textbf{Coverage never decreases on a touched file}, per this
    project's existing standard (\texttt{python-testing.md}, PL-TT-2:
    \texttt{pytest --co -q} before/after, count must not decrease) ---
    restated here because Conversation Mode's components (detectors, the
    state machine, the segmenter) are exactly the kind of new pure logic
    where it is cheapest to hold this line from the first commit, not
    retrofit it later.
  \item \textbf{The barge-in-vs-backchannel discrimination default
    (500\,ms, \fr{22}) ships with its calibration methodology as a
    reusable Tier 3 harness}, not just the number. \textbf{Enforcement:}
    this harness is one of the named entries required under
    \texttt{tests/conversation\_mode/harnesses/} by item 3 above, not a
    separate, unenforced requirement. The number came from one
    operator's live-tuned preference in one room; the harness that
    produced it (play a real passage, stop on a sustained-duration
    threshold, ask for feedback, adjust) is what makes that default
    re-tunable for a different voice or room later, per
    \S\ref{sec:oq}'s open question on this --- ship the method, not just
    its output.
\end{enumerate}

\end{document}
